% !TEX root = ../main.tex
\section{Theorie}
\label{sec:theorie}

\subsection{Wirk- und Blindleistung}

Die Kopplungsinduktivität verbindet zwei Spannungsquellen: den Wechselrichter mit der Grundschwingung $\underline{U}_{\mathrm{WR1}}$ und das Netz mit $\underline{U}_{\mathrm N}$. Die Spannung $\underline{U}_{\mathrm N}$ wird als Bezugszeiger gewählt, $\underline{U}_{\mathrm{WR1}}$ eilt um den Winkel $\Th$ nach. Bei vernachlässigtem Widerstand fliesst der Strom
\begin{equation}
  \underline{I} = \frac{\underline{U}_{\mathrm N}-\underline{U}_{\mathrm{WR1}}}{\mathrm{j}X},
  \qquad \underline{U}_{\mathrm{WR1}} = U_{\mathrm{WR1}}\,\mathrm{e}^{-\mathrm{j}\Th},
  \qquad X=\omega L .
\end{equation}
Mit $\underline{S}=\underline{U}_{\mathrm N}\,\underline{I}^{*}=P+\mathrm{j}Q$ (Bezugssinn aus Sicht des Netzes, passend zur Konvention in Abschnitt~\ref{sec:vorzeichen}) folgt
\begin{equation}
  P=\frac{\UN\,U_{\mathrm{WR1}}\,\sin\Th}{X},
  \qquad
  Q=\frac{\UN\,\bigl(\UN-U_{\mathrm{WR1}}\cos\Th\bigr)}{X}.
  \label{eq:PQ}
\end{equation}
Die Wirkleistung wird also hauptsächlich über den Winkel $\Th$ eingestellt, die Blindleistung hauptsächlich über die Amplitude $U_{\mathrm{WR1}}$ und damit über $\Udc$. Das Verhalten entspricht dem Zeigerdiagramm einer Synchronmaschine am Netz, wobei $U_{\mathrm{WR1}}$ der Polradspannung entspricht. Die Blindleistung der Grundschwingung wird null für $U_{\mathrm{WR1}}\cos\Th=\UN$.

\subsection{Spannungs- und Stromoberschwingungen bei Grundfrequenztaktung}

Bei Grundfrequenztaktung ist $\uWR$ eine Rechteckspannung mit den Amplituden $\pm\Udc$. Die Fourier-Zerlegung enthält nur ungerade Ordnungen $n$:
\begin{equation}
  \uWR(t)=\frac{4\,\Udc}{\pi}\sum_{n=1,3,5,\dots}\frac{1}{n}\sin(n\omega t),
  \qquad
  U_{\mathrm{WR},n}=\frac{4}{\pi\sqrt2}\,\frac{\Udc}{n}\approx 0.9\,\frac{\Udc}{n}.
  \label{eq:rechteck}
\end{equation}
Die Netzspannung wird als sinusförmig angenommen, deshalb treiben nur die Oberschwingungen von $\uWR$ Oberschwingungsströme durch die Induktivität:
\begin{equation}
  I_n=\frac{U_{\mathrm{WR},n}}{n\,X}.
  \label{eq:In}
\end{equation}
Die Stromoberschwingungen fallen deshalb mit $1/n^2$ ab. Mit $L=\SI{70}{\milli\henry}$ ist $X=2\pi\cdot\SI{50}{\hertz}\cdot\SI{70}{\milli\henry}=\SI{22.0}{\ohm}$. Für $\Udc=\SI{259}{\volt}$ ergeben sich $U_{\mathrm{WR},1}=\SI{233.1}{\volt}$ und $U_{\mathrm{WR},3}=\SI{77.7}{\volt}$, $U_{\mathrm{WR},5}=\SI{46.6}{\volt}$, $U_{\mathrm{WR},7}=\SI{33.3}{\volt}$ sowie $I_3=\SI{1.18}{\ampere}$, $I_5=\SI{0.42}{\ampere}$, $I_7=\SI{0.22}{\ampere}$ (Vergleich mit der Messung in Tabelle~\ref{tab:611}).

Ist $U_{\mathrm{WR1}}=\UN$ und $\Th=0$, so verschwindet der Grundschwingungsstrom ($I_1\approx0$), und der Strom besteht nur aus Oberschwingungen.

\subsection{Leistungsgrössen mit Oberschwingungen}

Die Wirkleistung $P$ ist der Mittelwert von $u\cdot i$ über alle Frequenzen, $P_1=U_1 I_1\cos\varphi_1$ gehört zur Grundschwingung (Index 1). Da die Netzspannung nahezu sinusförmig ist, gilt $P\approx P_1$. Der PM\,3000 bildet $Q=\sqrt{S^2-P^2}$; diese Grösse enthält die Verzerrungsblindleistung und wird deshalb nie null. Erst die Grundschwingungsblindleistung $Q_1=\pm\sqrt{(U_1I_1)^2-P_1^2}$ kann durch null gehen. Der Phasenwinkel der 1.\ Harmonischen ist am PM\,3000 immer $\SI{0}{\degree}$, weil er auf die eigene Grundschwingung bezogen ist, und eignet sich deshalb nicht zur Bestimmung von $Q_1$. Bei festem $P$ liegt $Q_1=0$ dort, wo $I_1$ minimal wird. Der Leistungsfaktor $\lambda=P/S$ ist wegen der Oberschwingungen kleiner als $\cos\varphi_1$.

\subsection{Pulsmuster}

Bei Pulsmustern mit eliminierten Oberschwingungen werden Kerben an berechneten Winkeln gesetzt, die gezielt einzelne Ordnungen aus $\uWR$ entfernen. Beim Trägerverfahren folgt die Pulsbreite einem sinusförmigen Sollwert, die Oberschwingungen liegen dann um die Schaltfrequenz. Beide Verfahren verkleinern die Grundschwingung $U_{\mathrm{WR1}}$ pro Volt $\Udc$ gegenüber dem Rechteck; für denselben Lastpunkt ist deshalb eine höhere Zwischenkreisspannung nötig. Eine höhere Schaltfrequenz verschiebt die Oberschwingungen zu höheren Ordnungen, wo die Induktivität stärker dämpft, bedeutet aber mehr Schaltverluste und mehr EMV-Störungen (qualitative Aussage, nicht gemessen).
